\begin{dang}{Điều kiện giao thoa sóng}
\Anlythuyet{
\circEX{1} \textbf{Kiểm tra tại M là cực đại hay cực tiểu}
\begin{itemize}
\item \textit{Phương pháp:} Tính $m =\dfrac{d_1-d_2}{\lambda}$
\begin{itemize}
\item Nếu m nguyên thì tại M là cực đại giao thoa.
\item Nếu m bán nguyên thì tại M là cực tiểu giao thoa.
\end{itemize}
\end{itemize}
\circEX{2} \textbf{Hai vân cùng loại đi qua hai điểm}

Giả sử hai vân cùng loại bậc k và bậc ($k+a$) đi qua hai điểm M và M’ thì
$\begin{cases}
MA-MB=k\lambda.\\
M'A-M'B=(k+a)\lambda.\\
\end{cases}$
\Note{
\begin{itemize}
\item Khoảng cách giữa hai cực đại hoặc hai cực tiểu liên tiếp trên đoạn giữa hai nguồn là $\dfrac{\lambda}2$.
\item Khoảng cách giữa một cực đại và một cực tiểu gần nhất trên đoạn giữa hai nguồn là $\dfrac{\lambda}4$.
\end{itemize}
}
}
\end{dang}
